\section{Defensa ante contenido conocido: capacidades y límites del Hash Matching}

Cuando una institución confiable ya verificó un archivo y este se incorporó conforme a la ley al reference set, la defensa escalable más madura es el matching. Evita que el modelo o una persona tengan que volver a interpretar el mismo contenido cada vez y es adecuado para atender la propagación repetida a gran escala; sin embargo, solo puede identificar contenido «lo bastante idéntico a un objeto conocido», no puede descubrir todo el contenido nuevo y tampoco puede sustituir la evaluación de legalidad ni el proceso de reporte de la plataforma.

\subsection{Cryptographic hash y perceptual hash}

\term{cryptographic hash} (hash criptográfico) asigna a una secuencia de bytes un resumen de longitud fija y es adecuado para determinar si un archivo es idéntico byte por byte; una compresión o un recorte leves cambian el resumen. \term{perceptual hash} (hash perceptual) extrae la estructura visual o auditiva, de modo que medios similares que pasaron por transformaciones comunes todavía pueden quedar en representaciones cercanas; mejora la recuperación de variantes, pero requiere umbrales, reference provenance y control de las coincidencias erróneas.

\lecturefigure{03-hash-matching.png}{El Hash matching es la defensa mínima para bloquear a gran escala material conocido y verificado.}{Subtítulos locales 05:10--06:00; mecanismo oficial de Thorn Safer Match; redibujo conceptual.}

\begin{knowledgebox}{Lectura de la figura: Match no es «volver a identificar el contenido»}
La plataforma primero calcula el hash o el fingerprint del incoming media y luego lo compara con un reference set controlado; si hay coincidencia, el caso pasa a la ruta de policy, revisión o reporte de la plataforma. Un cryptographic match suele aportar evidencia más sólida de byte identity, mientras que un perceptual match expresa similitud. En ambos casos hay que registrar hash algorithm/version, reference source, threshold, match confidence y las acciones humanas posteriores, para poder reproducir y auditar los resultados después de actualizar el sistema.
\end{knowledgebox}

\begin{warningbox}{Una Hash database no es una blocklist común}
Los Reference hash están asociados a evidencia de sensibilidad extremadamente alta, por lo que deben restringirse las fuentes de importación, los permisos de acceso, la capacidad de exportación y la visibilidad de los registros. La plataforma no puede ignorar la membership inference, el abuso interno, el etiquetado erróneo, la filtración entre tenants ni las obligaciones de eliminación por el hecho de que «solo guarda hashes». El servicio de matching debe devolver el resultado mínimo necesario, en lugar de exponer el contenido de las reference o una interfaz enumerable.
\end{warningbox}

\subsection{Por qué la coincidencia con contenido conocido aún requiere intervención humana y políticas}

Esta sección reconecta la señal de hash con la cadena de respuesta. Un mismo archivo puede aparecer en la obtención de evidencia para una denuncia, en reportajes periodísticos, en investigación, en sistemas legítimos de las autoridades o en el entorno de carga de un usuario, y en cada caso el contexto y las obligaciones legales pueden ser distintos; las distintas jurisdicciones, tipos de servicio y funciones de producto también tienen requisitos de reporte diferentes. Por ello, el equipo de ingeniería debe modelar el «match» como una señal de alta confianza, y no como una orden universal de eliminación que se salte toda policy.

\begin{importantbox}{Campos mínimos de un evento de coincidencia auditable}
Un durable match event contiene como mínimo: content identifier, hash algorithm/version, reference list/version, match distance, tenant/account, event time, policy version, review/action, reporting handoff y retry state. Los campos sensibles deben almacenarse en capas según el principio de mínimo privilegio, y deben diseñarse procesos explícitos para la eliminación, el legal hold y la incident response.
\end{importantbox}

\subsection{Resumen de la sección}

El Hash matching obtiene velocidad y consistencia a cambio de depender de un reference set verificado, y es la base de la known-material defense. Sus límites son igual de claros: no puede cubrir contenido desconocido, no puede inferir la identidad ni la intención de quien actúa, ni exime de la revisión del contexto ni de los procedimientos legales.
